% ============================================================
% SEÇÃO 10 – CONCLUSÕES
% ============================================================

\section{CONCLUSÕES}

O desenvolvimento do projeto LifeCity ao longo dos dois semestres possibilitou
uma imersão profunda em um campo tecnológico em plena expansão, revelando
oportunidades concretas para a criação de soluções digitais voltadas à comunicação
entre cidadãos e poder público. A construção deste aplicativo permitiu identificar
lacunas no mercado e propor um sistema com potencial de escalabilidade, capaz de
evoluir e atender demandas reais da sociedade.

No primeiro semestre, o projeto foi estruturado com foco na validação da ideia e
no desenvolvimento das funcionalidades essenciais: mapa interativo com pins de
eventos e reclamações, sistema de publicação de denúncias com geolocalização,
cadastro e autenticação de usuários, e gerenciamento básico de perfil. A
prototipação no Figma e a pesquisa com usuários confirmaram a relevância da
solução e orientaram as decisões de design e escopo.

No segundo semestre, o projeto avançou em maturidade técnica e experiência do
usuário. O backend foi migrado para uma solução própria em Node.js/Express com
autenticação JWT robusta, incluindo access token de curta duração, refresh token de
longa duração e renovação automática de sessão. O aplicativo passou por um
redesign visual completo, adotando modo escuro como padrão com suporte a
alternância entre temas claro e escuro. Foram incorporadas a tela de configurações
com gestão de conta, a identidade visual com o mascote do LifeCity e melhorias
significativas na usabilidade geral. Adicionalmente, foram implementadas
funcionalidades sociais como curtidas, comentários, sistema de amizades, aba de
destaques por engajamento e visualização de perfis públicos.

No terceiro semestre, o foco foi a camada de gamificação e colaboração. Foi
implementado o sistema de XP e Níveis, que recompensa o usuário por ações cívicas
como criar reclamações, receber curtidas e comentários. Sobre essa base, foi
construído o sistema de Conquistas, que desbloqueia automaticamente recompensas
ao atingir marcos predefinidos e exibe até três conquistas em destaque no perfil
público. Um centro de Notificações foi integrado para comunicar ao usuário curtidas,
comentários, pedidos de amizade, conquistas desbloqueadas e eventos de missão.

As Missões diárias e semanais introduziram objetivos estruturados: a cada novo
período, o sistema sorteia automaticamente uma missão do pool de templates e
acompanha o progresso em tempo real conforme reclamações são registradas. O
sistema de Equipes permite que grupos de até sete moradores colaborem em missões
semanais, recebendo bônus de XP proporcional ao desempenho coletivo — gerando
um incentivo à cooperação e à regularidade de uso do aplicativo.

Ao longo do processo, diversos conhecimentos adquiridos no curso de Engenharia
de Software foram aplicados de maneira prática e integrada: metodologia ágil Scrum
para organização das entregas, Análise de Projeto para estruturar requisitos e casos
de uso, Verificação e Validação para os casos de teste, Interação Humano-
Computador para o design das interfaces, e tecnologias de programação para web e
dispositivos móveis para a construção da solução.

Assim, o LifeCity representa não apenas um produto digital funcional, mas também
a materialização do aprendizado acumulado ao longo da formação acadêmica,
demonstrando a capacidade de integrar teoria, prática e inovação para gerar impacto
social real. A arquitetura incremental adotada — partindo do MVP com mapa e
reclamações, evoluindo para o ecossistema de gamificação e colaboração — mostrou-se
eficaz para entregar valor a cada fase sem comprometer a estabilidade das
funcionalidades anteriores.
